\documentclass[11pt,a4paper]{article}

\usepackage[utf8]{inputenc}
\usepackage[T1]{fontenc}
\usepackage{lmodern}
\usepackage[margin=2.4cm]{geometry}
\usepackage{microtype}
\usepackage{parskip}
\usepackage{booktabs}
\usepackage{longtable}
\usepackage{amssymb}
\usepackage{enumitem}
\usepackage{xcolor}
\usepackage{listings}
\usepackage{hyperref}
\usepackage{titlesec}
\usepackage{fancyhdr}

\definecolor{qahblue}{HTML}{1F4E79}
\definecolor{qahgreen}{HTML}{1E7B45}
\definecolor{qahgray}{HTML}{4A4A4A}
\definecolor{codebg}{HTML}{F4F6F8}
\definecolor{codecomment}{HTML}{2E7D32}
\definecolor{codestring}{HTML}{A31515}

\hypersetup{
  colorlinks=true,
  linkcolor=qahblue,
  urlcolor=qahblue,
  citecolor=qahblue,
  pdftitle={Contributing a Plugin to the QAH Git Marketplace},
  pdfauthor={Acxel Orozco}
}

\titleformat{\section}{\Large\bfseries\color{qahblue}}{\thesection}{0.7em}{}
\titleformat{\subsection}{\large\bfseries\color{qahblue}}{\thesubsection}{0.6em}{}

\pagestyle{fancy}
\fancyhf{}
\fancyhead[L]{\small QAH Git Marketplace}
\fancyhead[R]{\small Contributor guide}
\fancyfoot[C]{\thepage}
\renewcommand{\headrulewidth}{0.4pt}

\lstdefinestyle{qah}{
  backgroundcolor=\color{codebg},
  basicstyle=\ttfamily\small,
  breaklines=true,
  breakatwhitespace=false,
  commentstyle=\color{codecomment},
  keywordstyle=\color{qahblue}\bfseries,
  stringstyle=\color{codestring},
  frame=single,
  rulecolor=\color{qahgray!40},
  columns=fullflexible,
  keepspaces=true,
  showstringspaces=false,
  numbers=left,
  numberstyle=\tiny\color{qahgray},
  numbersep=8pt,
  captionpos=b
}
\lstset{style=qah}

\newcommand{\code}[1]{\texttt{#1}}
\newcommand{\repo}{\href{https://github.com/ADOBApps/qah-git-marketplace}{ADOBApps/qah-git-marketplace}}
\newcommand{\pluginrepo}{\href{https://github.com/ADOBApps/qah-git-marketplace/releases/latest}{Git Plugin Marketplace}}

\title{\textbf{Contributing a Plugin to the\\QAH Git Marketplace}\\[0.4em]
\large A practical guide for external users: from a plugin folder to a merged pull request}
\author{Acxel Orozco\\Quantum Analysis Helper}
\date{November 2026 \quad -- \quad Git Plugin Marketplace 1.1.1}

\begin{document}
\maketitle
\thispagestyle{fancy}

\begin{abstract}
This document explains how any GitHub user (a ``custom user'') can ask to upload
a plugin to the public QAH Git Marketplace. Uploads are **pull requests**: a
contributor publishes the plugin to their own fork, opens a pull request, the
validation workflow checks the catalog and archives, and the repository owner
reviews and merges. Deletion is owner-only. No contributor ever receives the
owner's token.
\end{abstract}

\tableofcontents
\newpage

% ---------------------------------------------------------------------------
\section{Purpose and audience}

The QAH Git Marketplace is a public Git repository that stores a catalog
(\code{index.json}) and plugin archives (\code{.qplugin.zip}) for Quantum
Analysis Helper (QAH). It is read by the \emph{Git Plugin Marketplace} plugin.

\begin{description}[leftmargin=1.5em,style=nextline]
  \item[Upstream] \repo, branch \code{master}. Public read; protected branch.
  \item[Fork] your own copy, \code{https://github.com/<you>/qah-git-marketplace}.
  \item[Pull request (PR)] your proposed change, reviewed and merged by the owner.
  \item[Token] a personal access token created in \emph{your} GitHub account.
  \item[CI] the \code{validate} GitHub Actions check that must pass before merging.
\end{description}

There are two supported workflows:

\begin{enumerate}[leftmargin=1.5em]
  \item \textbf{Automatic} (recommended): the plugin or CLI forks the repository,
        commits to a branch and opens the PR for you.
  \item \textbf{Manual}: you fork and publish to your fork yourself, then open
        the PR on GitHub.
\end{enumerate}

Both produce the same result. The owner merges; the plugin becomes available to
everyone after users refresh the marketplace.

% ---------------------------------------------------------------------------
\section{Before you start}

\subsection{What you need}

\begin{itemize}[leftmargin=1.5em]
  \item A GitHub account.
  \item A completed QAH plugin folder (at minimum \code{manifest.json} and
        \code{\_\_init\_\_.py}).
  \item A \emph{new} version number for this submission.
  \item Your own GitHub token (Section~\ref{sec:token}).
  \item Git Plugin Marketplace 1.1.1 or newer, or the standalone
        \code{git\_marketplace\_cli.py}.
\end{itemize}

\subsection{Public source warning}

Everything committed to the public marketplace is publicly downloadable. A QAH
Python plugin archive normally contains its \code{.py} source files. Do not
submit private, paid, NDA-covered or secret code to the public repository. Use a
separate private staging repository for such plugins.

\subsection{Plugins must be distributable by you}

Submit only plugins you own or are allowed to redistribute. The owner may reject
a submission that appears to violate licensing or ownership.

% ---------------------------------------------------------------------------
\section{Plugin folder requirements}

The publisher accepts the classic QAH plugin layout:

\begin{lstlisting}[caption={Minimal plugin folder}]
my_plugin/
  manifest.json      required
  __init__.py        required (Plugin class)
  README.md          optional, copied to the catalog
  icon.svg           optional
  ...                any other runtime files
\end{lstlisting}

Required manifest fields:

\begin{center}
\begin{tabular}{@{}llp{7.2cm}@{}}
\toprule
Field & Example & Notes \\
\midrule
\code{id} & \code{my\_plugin} & letters, digits, \code{.}, \code{\_} or \code{-} \\
\code{name} & My Plugin & display name \\
\code{version} & \code{1.2.0} & new version; semantic versioning recommended \\
\code{author} & Your Name & non-empty \\
\code{description} & Short text & non-empty \\
\code{category} & \code{tools} & \code{analysis}, \code{development}, \code{tools}, \code{visualization} \\
\bottomrule
\end{tabular}
\end{center}

Example:

\begin{lstlisting}[language=Java,caption={manifest.json}]
{
  "name": "My Plugin",
  "id": "my_plugin",
  "version": "1.2.0",
  "author": "Your Name",
  "description": "What the plugin does",
  "type": "python",
  "category": "tools",
  "main": "__init__.py",
  "dependencies": [],
  "license": { "type": "free", "max_activations": 3 }
}
\end{lstlisting}

The publisher builds the archive with \code{manifest.json} at the root. Do not
ship a wrapper folder, \code{.git}, \code{\_\_pycache\_\_}, \code{dist},
\code{target}, virtual environments or old \code{.qplugin.zip} files; they are
excluded automatically.

% ---------------------------------------------------------------------------
\section{Create your own GitHub token}\label{sec:token}

A token is a personal credential. It is created in \emph{your} GitHub account;
the owner cannot create one for you and must never share theirs.

\subsection{Option 1: GitHub CLI (easiest)}

\begin{lstlisting}[language=bash,caption={Terminal}]
gh auth login
\end{lstlisting}

After that the plugin can read the token with \code{gh auth token}; no manual
copying is needed as long as \code{gh} is on \code{PATH}.

\subsection{Option 2: classic personal access token (recommended for PRs)}

\begin{enumerate}[leftmargin=1.5em]
  \item Sign in to your GitHub account.
  \item Open \url{https://github.com/settings/tokens/new}.
  \item Note: \code{QAH Git Marketplace contributor}.
  \item Expiration: 30, 60 or 90 days.
  \item Scope: \code{public\_repo}.
  \item Generate the token and copy it once.
\end{enumerate}

\code{public\_repo} can fork the public marketplace, write to your fork and open
pull requests. This is the simplest option for the automatic submission.

\subsection{Option 3: fine-grained token}

\begin{enumerate}[leftmargin=1.5em]
  \item Open \url{https://github.com/settings/personal-access-tokens/new}.
  \item Resource owner: your account (or the organisation that owns your fork).
  \item Repository access:
    \begin{itemize}
      \item upstream \code{ADOBApps/qah-git-marketplace}:
            \textbf{Contents: Read}, \textbf{Pull requests: Read and write};
      \item your fork \code{<you>/qah-git-marketplace}:
            \textbf{Contents: Read and write}.
    \end{itemize}
  \item Metadata: Read (mandatory).
  \item Generate and copy the token.
\end{enumerate}

\textbf{Fork caveat.} A fine-grained token can only grant access to repositories
that exist when the token is created. If you have not forked yet, either:
\begin{itemize}[leftmargin=1.5em]
  \item create the fork manually first, then create the token; or
  \item use a classic \code{public\_repo} token; or
  \item grant the token access to ``All repositories'' (not recommended).
\end{itemize}

\subsection{Token security}

\begin{itemize}[leftmargin=1.5em]
  \item Never send your token to the owner or paste it into a public issue.
  \item Never commit it. The plugin stores it in
        \code{\textasciitilde/.qah/git\_marketplace.json} (mode 0600 on
        Linux/macOS; user-profile ACL on Windows).
  \item Use the minimum scope and a short expiration; revoke it when done.
  \item A read-only token can browse and install but cannot publish or delete.
\end{itemize}

% ---------------------------------------------------------------------------
\section{Configure the plugin}

Open \textbf{QAH $\rightarrow$ Plugins $\rightarrow$ Git Plugin Marketplace
$\rightarrow$ Settings} and fill in:

\begin{center}
\begin{tabular}{@{}ll@{}}
\toprule
Field & Value \\
\midrule
Owner & \code{ADOBApps} \\
Repository & \code{qah-git-marketplace} \\
Branch & \code{master} \\
Token & your token (Section~\ref{sec:token}) \\
\bottomrule
\end{tabular}
\end{center}

Click \textbf{Test connection}. A read-only token reports \code{write: no}; that
is expected for contributors. The public catalog remains readable without any
token.

The same settings can be written from the CLI:

\begin{lstlisting}[language=bash,caption={CLI configuration}]
py -3 git_marketplace_cli.py config --owner ADOBApps \
    --repo-name qah-git-marketplace --branch master
py -3 git_marketplace_cli.py doctor
\end{lstlisting}

% ---------------------------------------------------------------------------
\section{Option A: automatic pull request (recommended)}

\subsection{GUI}

\begin{enumerate}[leftmargin=1.5em]
  \item Open \textbf{Git Plugin Marketplace}.
  \item Click \textbf{Publish Plugin...}.
  \item Select your plugin folder. The dialog reads \code{manifest.json}.
  \item Review/edit name, author, category, description and changelog.
  \item Tick \textbf{Submit as a pull request (fork + PR)}.
        It is pre-selected automatically when your token is read-only.
  \item Click \textbf{Submit pull request}.
\end{enumerate}

The plugin creates or reuses your fork, commits the catalog entry and archive to
a temporary branch, opens the PR and shows the URL.

\subsection{CLI}

\begin{lstlisting}[language=bash,caption={Automatic submission}]
py -3 git_marketplace_cli.py submit "C:\path\to\my_plugin" \
    --changelog "Initial release: what it does"
\end{lstlisting}

Useful flags: \code{--version 1.2.0}, \code{--draft}, \code{--title "..."},
\code{--homepage "..."}. Without \code{--version} the manifest version is used.

\subsection{What happens behind the scenes}

\begin{enumerate}[leftmargin=1.5em]
  \item The folder is validated and archived as a \code{.qplugin.zip}.
  \item The current upstream \code{index.json} is read and a new entry is added.
  \item Your GitHub login is identified from the token.
  \item Your fork is created (if missing) and synced with upstream.
  \item A branch \code{marketplace/<id>-<version>-<timestamp>} is created from
        the upstream head.
  \item \code{index.json}, \code{plugins/<id>/manifest.json}, optional README and
        the ZIP are committed to that branch in \textbf{one commit}.
  \item A pull request is opened against \code{ADOBApps:master}.
  \item If the PR cannot be opened, the temporary branch is deleted.
\end{enumerate}

For the repository owner, the same action uses a branch in the upstream
repository instead of a fork (owner same-repository PR).

% ---------------------------------------------------------------------------
\section{Option B: manual fork and pull request}

Use this if you prefer Git over the automatic action.

\begin{enumerate}[leftmargin=1.5em]
  \item On GitHub, open \repo{} and click \textbf{Fork}.
  \item Configure the plugin to point at your fork:
        Owner \code{<you>}, Repository \code{qah-git-marketplace},
        Branch \code{master}, your token.
  \item Publish to your fork:
\begin{lstlisting}[language=bash,caption={Publish to your fork}]
py -3 git_marketplace_cli.py --repo <you>/qah-git-marketplace \
    --branch master publish "C:\path\to\my_plugin" \
    --changelog "What changed"
\end{lstlisting}
  \item On GitHub, open a pull request from \code{<you>:master} to
        \code{ADOBApps:master}.
  \item Switch your plugin settings back to \code{ADOBApps} if you want to read
        the upstream catalog again.
\end{enumerate}

If the upstream moved while you were working, sync your fork (GitHub's
\textbf{Sync fork} button or the automatic action) and publish again. The PR
updates automatically.

% ---------------------------------------------------------------------------
\section{Validation: the \code{validate} check}

Every PR runs \code{.github/workflows/validate-marketplace.yml}. The job checks:

\begin{itemize}[leftmargin=1.5em]
  \item \code{index.json} is valid JSON with schema 1;
  \item plugin ids and versions are unique and safe;
  \item every local ZIP exists and its size and SHA256 match the catalog;
  \item the ZIP contains \code{manifest.json} at its root;
  \item the manifest \code{id} and \code{version} match the catalog entry;
  \item required manifest fields are non-empty;
  \item there are no absolute paths, \code{..} entries or symlinks;
  \item the catalog \code{version} equals the highest non-yanked version.
\end{itemize}

A green check means the data is consistent. It does not mean the plugin has been
reviewed; the owner still reviews and merges.

\subsection{Common validation failures}

\begin{center}
\begin{longtable}{@{}p{6.2cm}p{9.2cm}@{}}
\toprule
Message & Fix \\
\midrule
\code{sha256 mismatch} & Publish again; do not edit the ZIP by hand. \\
\code{size ... != catalog ...} & Rebuild and republish the same version. \\
\code{manifest.json is not at the archive root} & The publisher does this for you; do not repack with a wrapper folder. \\
\code{manifest id ... != ...} & The folder manifest and the catalog entry disagree; republish from the correct folder. \\
\code{manifest version ... != ...} & Use the same version in the manifest and \code{--version}. \\
\code{missing ...} & The catalog references a file that is not in the PR. \\
\code{catalog version ... != latest ...} & Publish a higher version or fix \code{index.json}. \\
\bottomrule
\end{longtable}
\end{center}

% ---------------------------------------------------------------------------
\section{Review, merge and publication}

\begin{itemize}[leftmargin=1.5em]
  \item \code{master} is protected: PR required, one approval, Code Owner
        review, required \code{validate} check, no force pushes, no deletions.
  \item Only the owner (\code{@ADOBApps}) can merge. Contributors cannot push to
        \code{master}.
  \item The owner may request changes; update your fork/PR and republish.
  \item After the merge, the catalog is live immediately. Users click
        \textbf{Refresh} in Git Plugin Marketplace and see the new version.
\end{itemize}

Please keep one plugin/version per PR and do not modify other plugins in the same
change.

% ---------------------------------------------------------------------------
\section{Updating an existing plugin}

Publish a new version with the same \code{id}. Older versions remain in the
catalog, so users can still install them.

\begin{lstlisting}[language=bash,caption={Publishing an update}]
py -3 git_marketplace_cli.py submit "C:\path\to\my_plugin" \
    --version 1.3.0 --changelog "Fix ..., add ..."
\end{lstlisting}

Never reuse a version number for different bytes. The catalog replaced the
entry if the version already exists, but this breaks reproducibility for users.

% ---------------------------------------------------------------------------
\section{Deletion and ownership}

\begin{itemize}[leftmargin=1.5em]
  \item Only the owner has a write token and can delete via the plugin
        (\textbf{Delete from marketplace}) or by merging a deletion PR.
  \item A contributor may open a PR that removes an entry, but it cannot be
        merged without the owner.
  \item Deletion removes the entry from the branch tip. Git history still
        contains the old objects; this is not a secure erase.
\end{itemize}

% ---------------------------------------------------------------------------
\section{Troubleshooting}

\begin{center}
\begin{longtable}{@{}p{6.0cm}p{9.4cm}@{}}
\toprule
Symptom & Cause and fix \\
\midrule
\code{401} / bad credentials & Token expired or mistyped. Create a new one and update Settings. \\
\code{403} on publish & Token lacks Contents write on the target repo, or the organisation requires SSO authorisation. \\
\code{403} on submit & Token cannot fork or open PRs. Use \code{public\_repo}, or grant Pull requests write on upstream. \\
\code{404} & Private repository without access, or the fork is not ready yet. Wait and retry; check the owner/repo spelling. \\
\code{422} on branch & Branch already exists (rare timestamp collision) or the branch moved. Retry. \\
\code{422} on PR & A PR for this branch may already exist, or the base branch changed. Check GitHub. \\
Fork creation timeout & GitHub is slow. Retry after a minute; the fork may already exist. \\
\code{writable: false} & Your token is read-only; use \textbf{Submit as a pull request}, not direct publish. \\
CI failed & Read the job log; usually a SHA/size mismatch, missing file, or version not highest. Republish the same version. \\
Plugin missing after merge & Click \textbf{Refresh}; restart QAH if the menu does not update. \\
\bottomrule
\end{longtable}
\end{center}

If the problem persists, open an issue at \repo{} with the command, the exact
error text and the plugin id/version. Never paste your token.

% ---------------------------------------------------------------------------
\appendix
\section{Catalog structure}

\begin{lstlisting}[language=Java,caption={index.json (abridged)}]
{
  "schema": 1,
  "name": "QAH Git Marketplace",
  "updated": "2026-11-06T22:31:00Z",
  "plugins": [
    {
      "id": "my_plugin",
      "name": "My Plugin",
      "version": "1.2.0",
      "author": "Your Name",
      "description": "What it does",
      "category": "tools",
      "dependencies": [],
      "versions": [
        {
          "version": "1.2.0",
          "path": "plugins/my_plugin/1.2.0/my_plugin-1.2.0.qplugin.zip",
          "sha256": "...64 hex characters...",
          "size": 12345,
          "changelog": "Initial release",
          "published": "2026-11-06T22:31:00Z",
          "yanked": false
        }
      ]
    }
  ]
}
\end{lstlisting}

\section{PR checklist}

\begin{itemize}[leftmargin=1.5em]
  \item[$\square$] I own or may redistribute this plugin.
  \item[$\square$] The folder has \code{manifest.json} and \code{\_\_init\_\_.py}.
  \item[$\square$] The version is new and matches the manifest.
  \item[$\square$] \code{validate} passes.
  \item[$\square$] No tokens, credentials or private data are included.
  \item[$\square$] I understand the public repository exposes the archive contents.
\end{itemize}

\section{Command reference}

\begin{lstlisting}[language=bash,caption={Git Plugin Marketplace CLI}]
# settings / diagnosis
git_marketplace_cli.py config --owner ADOBApps \
    --repo-name qah-git-marketplace --branch master --token <TOKEN>
git_marketplace_cli.py doctor

# automatic pull request
git_marketplace_cli.py submit <folder> --changelog "..." [--draft]

# manual fork workflow
git_marketplace_cli.py --repo <you>/qah-git-marketplace --branch master \
    publish <folder> --changelog "..."

# local validation / build
git_marketplace_cli.py validate <folder>
git_marketplace_cli.py build <folder> --out dist

# after the owner merges (public read, no token)
git_marketplace_cli.py list --refresh
git_marketplace_cli.py install <plugin-id>
\end{lstlisting}

\end{document}
